Chương này trình bày và phân tích toàn bộ kết quả mô phỏng từ ba thí nghiệm. Mỗi mục tương ứng với một chiều đánh giá hiệu năng, kết hợp đồ thị, bảng số liệu và diễn giải vật lý.

\subsection{Thí nghiệm 1: BLER vs $E_s/N_0$}
\label{sec:sim01}

\subsubsection{Điều kiện mô phỏng}

Kênh Rician K = 15 dB, LEO 600 km, SCS 30 kHz. MCS A (QPSK $r=1/2$). Dải $E_s/N_0$: $-12$ đến $+8$ dB. Số thử nghiệm: 20\,000.

\subsubsection{Kết quả}

\begin{figure}[H]
    \centering
    \includegraphics[width=\linewidth]{sim01_bler_ir_leo.png}
    \caption[BLER vs $E_s/N_0$ -- IR-HARQ, MCS A, LEO 600 km]{\bfseries\fontsize{12pt}{0pt}\selectfont
    BLER vs $E_s/N_0$ -- IR-HARQ, MCS A (QPSK r=1/2), K = 15 dB, LEO 600 km.
    Đường chấm xám: No HARQ baseline. Các đường xanh (đậm dần): TX = 1 đến TX = 4.
    Đường ngang đứt: mục tiêu BLER = $10^{-3}$.}
    \label{fig:sim01}
\end{figure}

Bảng~\ref{tab:sim01} tổng hợp điểm làm việc tại BLER = $10^{-3}$.

\begin{table}[H]
    \centering
    \caption[Điểm làm việc BLER = $10^{-3}$ -- IR-HARQ, MCS A]{\bfseries\fontsize{12pt}{0pt}\selectfont
    $E_s/N_0$ tại BLER = $10^{-3}$ -- IR-HARQ, MCS A (QPSK r=1/2), K = 15 dB, LEO 600 km}
    \label{tab:sim01}
    \begin{tabularx}{0.88\textwidth}{
    | >{\centering\arraybackslash}X
    | >{\centering\arraybackslash}X
    | >{\centering\arraybackslash}X |}
    \hline
    \bfseries Sơ đồ & \bfseries $E_s/N_0$ tại BLER=$10^{-3}$ & \bfseries Độ lợi vs No HARQ \\\hline
    No HARQ / IR TX=1 & $+2{,}0$ dB & -- \\\hline
    IR TX=2 & $-2{,}5$ dB & $4{,}5$ dB \\\hline
    IR TX=3 & $-5{,}0$ dB & $7{,}0$ dB \\\hline
    IR TX=4 & $-7{,}0$ dB & \textbf{9{,}0 dB} \\\hline
    \end{tabularx}
    \begin{flushleft}
    \small \textit{No HARQ và IR TX=1 cho kết quả giống nhau (cùng 1 lần phát). Ngưỡng AWGN lý thuyết cho $r=1/2$, $\Delta=1{,}5$ dB: $-2{,}3$ dB. Tại kênh Rician K=15 dB, BLER=$10^{-3}$ đạt ở $+2{,}0$ dB (lề fading Rician $\approx 4{,}3$ dB).}
    \end{flushleft}
\end{table}

\subsubsection{Phân tích}

\textbf{Độ lợi link rõ rệt.} Ở TX = 1, IR-HARQ và No HARQ cho BLER giống nhau (cùng 1 lần phát, chưa có tích lũy phân tập). Từ TX = 2 trở đi, mỗi lần phát bổ sung thêm thông tin tương hỗ mới theo phương trình (\ref{eq:ir_mi}), đường BLER dịch sang trái từ 2 đến 4,5 dB mỗi TX (lớn nhất ở TX=2: 4,5 dB, giảm dần về TX=4: 2,0 dB). Ở TX = 4, IR-HARQ đạt BLER = $10^{-3}$ tại $-7{,}0$ dB -- độ lợi \textbf{9{,}0 dB} so với No HARQ (cần $+2{,}0$ dB).

\textbf{Đặc tính waterfall dốc.} Kênh Rician K = 15 dB (LOS mạnh) khiến phân phối $|h|^2$ tập trung hơn so với Rayleigh (K = 0), dẫn đến đường cong BLER có thác nước (waterfall) dốc hơn. Khi tích lũy $n$ TX theo IR, bậc tự do của phân phối chi-bình phương tăng lên $2n$ nên phương sai tương đối giảm và đường cong waterfall ngày càng dốc -- phù hợp với lý thuyết xác suất và kết quả Hình 6 của~\cite{tuninato2025}.

\subsection{Thí nghiệm 2: Phân tích N$_{\min}$ theo SCS}
\label{sec:sim02}

\subsubsection{Kết quả}

\begin{figure}[H]
    \centering
    \includegraphics[width=0.92\linewidth]{sim03_nmin_leo.png}
    \caption[N$_{\min}$ theo SCS -- LEO 600 km và LEO 1200 km]{\bfseries\fontsize{12pt}{0pt}\selectfont
    Số tiến trình HARQ tối thiểu N$_{\min}$ để tránh pipeline stall theo SCS (LEO).
    Đường đứt đỏ: giới hạn N = 32 của TS 38.214~\cite{ts38214} Mục 5.1.
    Cột vượt giới hạn: không khả thi với chuẩn hiện tại.}
    \label{fig:sim02}
\end{figure}

\begin{table}[H]
    \centering
    \caption[Bảng N$_{\min}$ -- quỹ đạo LEO]{\bfseries\fontsize{12pt}{0pt}\selectfont
    N$_{\min}$ theo quỹ đạo LEO và SCS -- công thức (\ref{eq:nmin}), RTT từ TR 38.811~\cite{tr38811}}
    \label{tab:nmin_full}
    \begin{tabularx}{0.92\textwidth}{
    | >{\centering\arraybackslash}X
    | >{\centering\arraybackslash}X
    | >{\centering\arraybackslash}X
    | >{\centering\arraybackslash}X
    | >{\centering\arraybackslash}X |}
    \hline
    \bfseries Quỹ đạo & \bfseries SCS 15 kHz & \bfseries SCS 30 kHz & \bfseries SCS 60 kHz & \bfseries SCS 120 kHz \\\hline
    LEO 600 km  & \textbf{14} $\checkmark$ & \textbf{27} $\checkmark$ & \textcolor{red}{53} $\times$ & \textcolor{red}{105} $\times$ \\\hline
    LEO 1200 km & \textbf{26} $\checkmark$ & \textcolor{red}{50} $\times$ & \textcolor{red}{99} $\times$ & \textcolor{red}{196} $\times$ \\\hline
    \end{tabularx}
    \begin{flushleft}
    \small $\checkmark$: N$_{\min} \leq 32$ (khả thi không stall); $\times$: N$_{\min} > 32$ (pipeline stall không thể tránh).\\
    \textit{Giới hạn N = 32: TS 38.214~\cite{ts38214} Mục 5.1.}
    \end{flushleft}
\end{table}

\subsubsection{Phân tích}

Bảng~\ref{tab:nmin_full} và Hình~\ref{fig:sim02} cho thấy bức tranh định lượng rõ ràng về khả năng áp dụng HARQ ở LEO.

\textbf{Chỉ 3 trường hợp khả thi trong 8 tổ hợp LEO.} LEO 600 km ở SCS 15 kHz (N$_{\min}$ = 14) và SCS 30 kHz (N$_{\min}$ = 27) đều dưới giới hạn N = 32. LEO 1200 km chỉ khả thi ở SCS 15 kHz (N$_{\min}$ = 26). Đây là hệ quả trực tiếp của N$_{\min} \propto \text{RTT}/T_{\text{slot}}$.

\textbf{Nghịch lý SCS cao.} SCS lớn hơn rút ngắn $T_{\text{slot}}$ nhưng lại \textit{tăng} N$_{\min}$: SCS 120 kHz ở LEO 600 km cần tới \textbf{105 tiến trình}, gấp hơn 3 lần giới hạn. Nguyên nhân: băng thông lớn hơn không giúp rút ngắn RTT -- RTT là thời gian lan truyền vật lý, không phụ thuộc SCS. Đây là điểm mấu chốt được nhấn mạnh trong TR 38.821 Mục 6.4~\cite{tr38821}.

\textbf{Cấu hình được khuyến nghị cho LEO.} Với LEO 600 km, SCS 30 kHz và N = 32 là lựa chọn tối ưu: util = 100\% (không có pipeline stall), $T_{\text{slot}}$ ngắn đủ để đáp ứng yêu cầu độ trễ. LEO 1200 km phải dùng SCS 15 kHz mới tránh được stall.

\subsection{Thí nghiệm 3: Hiệu quả năng lượng trên bit}
\label{sec:sim03}

\subsubsection{Kết quả}

\begin{figure}[H]
    \centering
    \includegraphics[width=0.85\linewidth]{sim05_energy_per_bit.png}
    \caption[Năng lượng/bit vs SNR -- IR-HARQ, LEO 600 km]{\bfseries\fontsize{12pt}{0pt}\selectfont
    Năng lượng chuẩn hóa trên bit vs $E_s/N_0$ -- LEO 600 km, K = 15 dB.
    $P_{tx} = 1$ W, SCS 30 kHz, MCS A. So sánh IR-HARQ và No HARQ.}
    \label{fig:sim03}
\end{figure}

\begin{table}[H]
    \centering
    \caption[Năng lượng/bit tại các điểm SNR -- LEO 600 km]{\bfseries\fontsize{12pt}{0pt}\selectfont
    Năng lượng chuẩn hóa/bit -- LEO 600 km, K = 15 dB, $P_{tx} = 1$ W}
    \label{tab:sim03}
    \begin{tabularx}{0.92\textwidth}{
    | >{\centering\arraybackslash}X
    | >{\centering\arraybackslash}X
    | >{\centering\arraybackslash}X
    | >{\centering\arraybackslash}X
    | >{\centering\arraybackslash}X |}
    \hline
    \bfseries Sơ đồ & \bfseries $E_{\text{bit}}$ tại $-5$ dB & \bfseries $E_{\text{bit}}$ tại 0 dB & \bfseries $E_{\text{bit}}$ tại 5 dB & \bfseries avg\_ntx tại 5 dB \\\hline
    No HARQ & $\sim 0{,}54$ J/bit & $\sim 5{,}2\times 10^{-4}$ J/bit & $\sim 5{,}0\times 10^{-4}$ J/bit & 1{,}00 \\\hline
    IR-HARQ & $\sim 1{,}1\times 10^{-3}$ J/bit & $\sim 5{,}2\times 10^{-4}$ J/bit & $\sim 5{,}0\times 10^{-4}$ J/bit & $\sim 1{,}00$ \\\hline
    \end{tabularx}
\end{table}

\subsubsection{Phân tích}

\textbf{Độ lợi năng lượng vượt trội ở SNR thấp.} Hình~\ref{fig:sim03} và Bảng~\ref{tab:sim03} cho thấy ở $E_s/N_0 = -5$ dB, No HARQ cần $\sim 0{,}54$ J/bit vì BLER $\approx 99{,}9\%$ -- gần như toàn bộ năng lượng bị lãng phí cho các gói không thành công. IR-HARQ giảm xuống còn $\sim 1{,}1 \times 10^{-3}$ J/bit, tức cải thiện \textbf{$\sim 500$ lần (gần 3 bậc độ lớn)} nhờ tích lũy MI từ 4 TX đưa BLER cuối về gần 0 trong khi No HARQ hầu như không giao được bit nào.

\textbf{Hội tụ ở SNR cao.} Ở $E_s/N_0 > 5$ dB, cả hai sơ đồ hội tụ về cùng $E_{\text{bit}} \approx P_{tx} \cdot T_{\text{slot}} / k \approx 3\times 10^{-3}$ J/bit vì avg\_ntx $\to$ 1 và BLER $\to$ 0. Điều này cho thấy HARQ đặc biệt có giá trị ở vùng SNR hạn chế -- đúng với kịch bản thiết bị IoT vệ tinh chạy bằng pin.

\textbf{IR-HARQ với LEO 600 km -- cấu hình lý tưởng.} Vì util = 100\% (N = 32 $\geq$ N$_{\min}$ = 27 ở SCS 30 kHz), toàn bộ lợi ích năng lượng của IR-HARQ được khai thác mà không có chi phí pipeline stall. Đây là cấu hình tham chiếu cho thiết kế hệ thống NTN LEO thực tế.

\subsection{Tổng hợp}
\label{sec:summary}

\begin{table}[H]
    \centering
    \caption[Tổng hợp kết quả -- IR-HARQ, LEO 600 km]{\bfseries\fontsize{12pt}{0pt}\selectfont
    Tổng hợp kết quả đánh giá IR-HARQ tại LEO 600 km, K = 15 dB, MCS A, SCS 30 kHz}
    \label{tab:summary}
    \begin{tabularx}{0.98\textwidth}{
    | >{\centering\arraybackslash}X
    | >{\centering\arraybackslash}X
    | >{\centering\arraybackslash}X |}
    \hline
    \bfseries Chỉ tiêu & \bfseries No HARQ & \bfseries IR-HARQ (TX=4) \\\hline
    Độ lợi SNR tại BLER=$10^{-3}$ & 0 & \textbf{+9{,}0 dB} \\\hline
    $E_s/N_0$ tại BLER=$10^{-3}$ & $+2{,}0$ dB & $-7{,}0$ dB \\\hline
    Util (N=32, SCS 30 kHz) & 100\% & 100\% \\\hline
    $E_{\text{bit}}$ tại $-5$ dB & $\sim 0{,}54$ J/bit & $\sim 1{,}1\times 10^{-3}$ J/bit \\\hline
    Tỷ lệ cải thiện $E_{\text{bit}}$ & -- & \textbf{$\sim 500$ lần} \\\hline
    \end{tabularx}
\end{table}

\textbf{Nhận xét hệ thống:} Ba kết luận chính từ ba thí nghiệm:

\begin{enumerate}
    \item \textbf{IR-HARQ mang lại độ lợi link đáng kể:} 9{,}0 dB ở BLER = $10^{-3}$ so với No HARQ, nhờ tích lũy thông tin tương hỗ độc lập qua từng TX theo bất đẳng thức Jensen (\ref{eq:jensen}).
    \item \textbf{Pipeline stall là ràng buộc thiết kế quan trọng:} Chỉ LEO 600 km / SCS $\leq$ 30 kHz và LEO 1200 km / SCS 15 kHz đáp ứng N$_{\min} \leq 32$. SCS cao hơn nghịch lý làm tăng N$_{\min}$ vì RTT không thay đổi.
    \item \textbf{HARQ tiết kiệm $\sim 500$ lần năng lượng} so với No HARQ ở $E_s/N_0 = -5$ dB -- đặc biệt quan trọng cho IoT vệ tinh. Tại LEO 600 km với util = 100\%, IR-HARQ khai thác lợi ích này mà không mất thêm năng lượng do pipeline stall.
\end{enumerate}
